% GENERATED FILE. Edit canonical lessons, then run npm run generate:books.
% canonical-source-hash: fnv1a64:6ea3dd8545310989
% canonical-lessons: FR-C184-bouger, FR-C184-cueillir, FR-C184-defendre, FR-C184-disparaitre, FR-C184-emprunter

\chapter{To stir, To pick, To defend, To disappear, To borrow}
\label{ch:fr-to-stir-to-pick-to-defend-to-disappear-to-borrow}
\hlchaptermodality{184}

\begin{chapteropening}
\textbf{By the end of this chapter:} \emph{I can say to stir, to pick, to defend, to disappear and to borrow.}

Complete the last lesson of chapter 184: I can say to stir, to pick, to defend, to disappear and to borrow.
\end{chapteropening}

\section[bouger]{bouger — to stir, to move}

\label{lesson:FR-C184-bouger}



\begin{quote}
Before the new one: say the French for to ski, then the French for to slip.
\end{quote}

\subsection*{You'll want to know: bouger}
\textbf{bouger} — \textquotedblleft{}to stir, to move\textquotedblright{}.

\begin{grammarlens}[title={What you've built}]
One more word for this chapter.
\end{grammarlens}

\subsection*{Your turn}
\begin{itemize}
\raggedright
  \item \emph{Say it:} \emph{bouger}
  \item \emph{Say it:} \emph{bouger}, once more
  \item \emph{Say it:} say \emph{bouger}
  \item \emph{Recall:} say the French for to ski, then the French for to slip, then say \emph{bouger} again
\end{itemize}

\begin{tcolorbox}[breakable,colback=teal!4,colframe=teal!35!black,title={Before you move on}]
What does bouger mean? (\textquotedblleft{}To stir, to move\textquotedblright{}.) Say it once more.
\end{tcolorbox}
\section[cueillir]{cueillir — to pick (flowers, fruit)}

\label{lesson:FR-C184-cueillir}



\begin{quote}
Before the new one: say the French for to slip, then the French for to stir.
\end{quote}

\subsection*{You'll want to know: cueillir}
\textbf{cueillir} — \textquotedblleft{}to pick (flowers, fruit)\textquotedblright{}.

\begin{grammarlens}[title={What you've built}]
One more word for this chapter.
\end{grammarlens}

\subsection*{Your turn}
\begin{itemize}
\raggedright
  \item \emph{Say it:} \emph{cueillir}
  \item \emph{Say it:} \emph{cueillir}, once more
  \item \emph{Say it:} say \emph{cueillir}
  \item \emph{Recall:} say the French for to slip, then the French for to stir, then say \emph{cueillir} again
\end{itemize}

\begin{tcolorbox}[breakable,colback=teal!4,colframe=teal!35!black,title={Before you move on}]
What does cueillir mean? (\textquotedblleft{}To pick (flowers, fruit)\textquotedblright{}.) Say it once more.
\end{tcolorbox}
\section[défendre]{défendre — to defend}

\label{lesson:FR-C184-defendre}



\begin{quote}
Before the new one: say the French for to stir, then the French for to pick.
\end{quote}

\subsection*{You'll want to know: défendre}
\textbf{défendre} — \textquotedblleft{}to defend\textquotedblright{}.

\begin{grammarlens}[title={What you've built}]
One more word for this chapter.
\end{grammarlens}

\subsection*{Your turn}
\begin{itemize}
\raggedright
  \item \emph{Say it:} \emph{défendre}
  \item \emph{Say it:} \emph{défendre}, once more
  \item \emph{Say it:} say \emph{défendre}
  \item \emph{Recall:} say the French for to stir, then the French for to pick, then say \emph{défendre} again
\end{itemize}

\begin{tcolorbox}[breakable,colback=teal!4,colframe=teal!35!black,title={Before you move on}]
What does défendre mean? (\textquotedblleft{}To defend\textquotedblright{}.) Say it once more.
\end{tcolorbox}
\section[disparaître]{disparaître — to disappear}

\label{lesson:FR-C184-disparaitre}



\begin{quote}
Before the new one: say the French for to pick, then the French for to defend.
\end{quote}

\subsection*{You'll want to know: disparaître}
\textbf{disparaître} — \textquotedblleft{}to disappear\textquotedblright{}.

\begin{grammarlens}[title={What you've built}]
One more word for this chapter.
\end{grammarlens}

\subsection*{Your turn}
\begin{itemize}
\raggedright
  \item \emph{Say it:} \emph{disparaître}
  \item \emph{Say it:} \emph{disparaître}, once more
  \item \emph{Say it:} say \emph{disparaître}
  \item \emph{Recall:} say the French for to pick, then the French for to defend, then say \emph{disparaître} again
\end{itemize}

\begin{tcolorbox}[breakable,colback=teal!4,colframe=teal!35!black,title={Before you move on}]
What does disparaître mean? (\textquotedblleft{}To disappear\textquotedblright{}.) Say it once more.
\end{tcolorbox}
\section[emprunter]{emprunter — to borrow}

\label{lesson:FR-C184-emprunter}



\begin{quote}
Before the new one: say the French for to defend, then the French for to disappear.
\end{quote}

\subsection*{You'll want to know: emprunter}
\textbf{emprunter} — \textquotedblleft{}to borrow\textquotedblright{}.

\begin{grammarlens}[title={What you've built}]
One more word for this chapter.
\end{grammarlens}

\subsection*{Your turn}
\begin{itemize}
\raggedright
  \item \emph{Say it:} \emph{emprunter}
  \item \emph{Say it:} \emph{emprunter}, once more
  \item \emph{Say it:} say \emph{emprunter}
  \item \emph{Recall:} say the French for to defend, then the French for to disappear, then say \emph{emprunter} again
\end{itemize}

\begin{tcolorbox}[breakable,colback=teal!4,colframe=teal!35!black,title={Before you move on}]
What does emprunter mean? (\textquotedblleft{}To borrow\textquotedblright{}.) Say it once more.
\end{tcolorbox}
